\documentclass[10pt]{article}
\usepackage[margin=20mm]{geometry}
\usepackage{amsmath,amssymb,mathtools,fontspec,unicode-math,xurl,hyperref}
\IfFontExistsTF{DejaVu Serif}{\setmainfont{DejaVu Serif}}{\setmainfont{Latin Modern Roman}}
\setmathfont{latinmodern-math.otf}
\setlength{\parindent}{0pt}\setlength{\parskip}{6pt}\sloppy
\title{Initial scientific review: M25-0038}
\author{}
\date{25 October 2026}
\begin{document}
\maketitle
\section*{Identity, exact version, and outcome}
Ben Heuer and Daxin Xu, p-Adic non-abelian Hodge theory for curves via moduli stacks, DOI 10.1090/jams/1065. This review concerns the supplied 49-page author manuscript arXiv:2402.01365v1, submitted 2 February 2024, and its matching 3212-line `main.tex`. Page references below are physical PDF pages, numbered from the first page, and agree with the printed article numbering. PDF SHA-256:\par
> a20c65048e39da7e8b8ceee24570e1f0b1f
> b8b2feafa1814da45f8c22f3bff9c\par
No journal-final comparison was performed. None of the findings is represented as a verified defect of the journal version.\par
Outcome: Major, with two individually Minor findings. The proof of perfectoid GAGA invokes a false density statement (Problem 1). This is a demonstrated invalid premise in a chain supporting one of the paper’s main results. An exact local replacement argument is supplied and checked below; the GAGA theorem is not refuted. An introductory comparison isomorphism is false on a nonreduced spectral fibre, but the body explicitly uses the weaker correct comparison (Problem 2). A spectral-locus assertion lacks a genus restriction already present in later statements (Problem 3).\par
All 49 pages were read as extracted text and all 49 rendered source-page images were actually inspected. This is a systematic claim-focused review, not an independent proof of every theorem or cited dependency. The original TeX was searched throughout and selected evidence ranges were read; the complete TeX was not independently semantically audited.\par
\section*{Definitions, assumptions, and checked reasoning}
Throughout the paper, \(K\) is an algebraically closed complete extension of \(\mathbb Q_p\), \(X/K\) is a smooth projective connected curve, and \(G\) is a connected reductive group unless a narrower scope is specified. Write \(\Omega=\widetilde\Omega_X=\omega_X(-1)\). For \(G=\operatorname{GL}_n\), the Hitchin base is \[\mathcal A=\prod_{i=1}^n H^0(X,\Omega^{\otimes i})\otimes\mathbf G_a,\] and the spectral cover \(\pi:Z_b\to X\) is defined by the monic characteristic polynomial. The regular centralizer has Lie algebra \(\pi_*\mathcal O_{Z_b}\). In the discussion of relative cohomology, distinguish the projection \(q:X\times\mathcal A\to\mathcal A\) from \(\pi':Z\to\mathcal A\).\par
The following checks were made.\par
1. Sections 2–3, pp. 6–12: right torsor conventions, anti-diagonal contracted products and the inverse difference torsor; the logarithm’s kernel and the need for a \(p\)-surjective open subgroup for surjectivity; the distinction between left exactness in general and the stronger assertions subsequently proved for the algebraic centralizer.\par
2. Section 4, pp. 13–16: the use of commuting local representations in the Baker–Campbell–Hausdorff calculation, invariance of the canonical Higgs field under its exponential action, change-of-chart reduction to a faithful linear representation, and zero-Higgs descent conditional on the cited local Simpson correspondence.\par
3. Sections 5–6, pp. 16–29: regular-centralizer equivariance, the identity \(d a(\tau)=\varphi\), the sign in the local exponential twist, and cancellation by the inverse torsor. The constructed Higgs field on the resulting etale bundle is distinguished from the canonical Higgs field of that underlying etale bundle. The relative Hodge–Tate sequence, its vector-bundle base changes and the construction of the torsor of splittings were traced, conditional on the quoted cohomological input. Functoriality is restricted by the stated compatibility with Kostant sections.\par
4. Section 7, pp. 30–36: the comparison diagram and its kernel isomorphism, rather than an isomorphism of all Lie-algebra cohomology; the use of one-dimensional fibres; torsion freeness of \(\Lambda\) and the inverse-limit hypotheses; the restriction \(p\nmid |W|\) where Weyl invariants are used. The \(\operatorname{GL}_n\) result and the more restricted general reductive result were kept separate. The pointwise continuous exponential yields the claimed condensed comparison and homeomorphism, not an asserted analytic isomorphism.\par
5. Section 8, pp. 36–44: the double-coset argument, generation and finite presentation for vector bundles, the Fitting-ideal step, reduction from a curve to \(\mathbb P^1\) by a finite flat map, and the dependence of the Higgs-stack comparison on GAGA. The regular spectral-fibre/Picard argument was traced conditional on the cited representability and BNR input.\par
6. Section 9, pp. 44–48: representation stacks include a \(G\)-torsor on the test space; naive conjugacy classes alone do not give the required stack. The adjunction and faithful-representation reductions were checked. For the degree-zero argument, determinant and tensor reduction, Riemann–Hurwitz, and the growth of sections of a positive-degree subbundle on finite etale covers were checked: \(g'-1=n(g-1)\) and Riemann–Roch gives \(h^0\geq n(d-g+1)\), contradicting the bounded invariant-section space when the tensor power makes \(d>g-1\).\par
\section*{Problem 1: a false density premise in the GAGA proof}
\subsection*{Claim and exact location}
Lemma 8.2.7, p. 39, with proof in `main.tex`, lines 2398–2414, asserts that a double coset has a representative \[A'\in \operatorname{GL}_r(R\langle T^{\pm1}\rangle)
       \cap T^{-m}M_r(R\langle T\rangle).\] The proof claims that \(\operatorname{GL}_r(R[T^{\pm1}])\) is dense in \(\operatorname{GL}_r(R\langle T^{\pm1}\rangle)\), and uses this to obtain \[\operatorname{GL}_r(R\langle T^{\pm1}\rangle)
 =\operatorname{GL}_r(R[T^{\pm1}])U(R\langle T^{\pm1}\rangle),
 \qquad U(S)=1+\mathfrak m_KM_r(S^+).\] The claimed density, at source line 2404, is false.\par
\subsection*{Complete countercheck, including hypotheses}
Take \(p=3\), \(K=R=\mathbb C_3\), \(R^+=\mathcal O_{\mathbb C_3}\), and \(r=1\). Then \(\operatorname{Spa}(R,R^+)\) is an allowed affinoid perfectoid \(K\)-space. In the Gauss norm on \(D=\mathbb C_3\langle T^{\pm1}\rangle\), the element \[A=1+3T\] is a unit, with convergent inverse \(\sum_{j\geq0}(-3T)^j\).\par
Every unit of \(\mathbb C_3[T,T^{-1}]\) is \(cT^n\) for \(c\ne0\) and \(n\in\mathbb Z\). Indeed, the highest and lowest exponents of a nonzero Laurent polynomial add under multiplication over a field, so a product equal to \(1\) forces both factors to have exponent width zero. If \(n=0\), \[\|1+3T-c\|=\max\{|1-c|_3,1/3\}\geq 1/3.\] If \(n\ne0\), the constant coefficient of \(1+3T-cT^n\) is \(1\), so its norm is at least \(1\). Thus no Laurent-polynomial unit approximates \(A\) to distance less than \(1/3\). The permitted rank-one case already disproves density.\par
This does not disprove the displayed product decomposition: the particular \(A\) above belongs to \(U(D)\) and therefore has the factorization \(1\cdot A\). Nor does it disprove the conclusion of Lemma 8.2.7. The invalid step is precisely the proposed density justification.\par
\subsection*{Exact local repair}
Keep the preceding Lemma 8.2.5, which gives the required factorization inside \(U\). Ordinary Laurent-polynomial matrices, without requiring their inverses to be Laurent-polynomial, are dense in \(M_r(D)\). Given \(A\in\operatorname{GL}_r(D)\), choose \[P\in M_r(R[T^{\pm1}])\] sufficiently close to \(A\) that \(P\in\operatorname{GL}_r(D)\) and \(A_1=P^{-1}A\in U(D)\). This uses openness of invertibility and continuity of inversion, together with the open subgroup \(U\) already used in the paper. It does not claim \(P\in\operatorname{GL}_r(R[T^{\pm1}])\).\par
By the same vanishing/Čech argument as in the original proof, \(A_1=A_+A_-\) with \(A_+\in U(R\langle T\rangle)\) and \(A_-\in U(R\langle T^{-1}\rangle)\). Right multiplication by \(A_-^{-1}\) gives the representative \[A'=PA_+.\] It is analytically invertible. Since \(P\) has only finitely many negative powers and \(A_+\) has none, \(A'\in T^{-m}M_r(R\langle T\rangle)\) for sufficiently large \(m\). This proves exactly the lemma’s stated conclusion, conditional only on the preceding factorization input. The stronger Laurent-unit factorization need not be asserted.\par
\subsection*{Counterevidence, dependence, and impact}
The intended approximation mechanism is close to a valid one; the distinction is which ring contains the inverse of the approximating matrix. The repaired lemma feeds directly into the generation lemma 8.2.9, the finite-presentation argument 8.2.11 and the \(\mathbb P^1\) GAGA theorem 8.2.1. Those support the curve GAGA theorem 8.1.1 and the algebraic/analytic Higgs-stack comparison 8.4.6, announced as Theorem 1.7.1, as well as the regular Hitchin-fibre geometry.\par
Under the operative 0.7.1 calibration, an invalid derivation on this core dependency chain supports Major, even though an exact local repair is available. The evidence does not establish that any of these theorem conclusions is false. Focused AI acceptance confirms the stated local finding and repair; author adoption is unclaimed and human adjudication remains pending.\par
\section*{Problem 2: an introductory isomorphism contradicted by a spectral fibre}
\subsection*{Claim, location, and strongest charitable reading}
Theorem 1.4.1, p. 4, source lines 488–490, calls \[\psi:R^1 f_{v*}\operatorname{Lie}J\longrightarrow\Lambda\otimes_{\mathbb Z_p}\mathbf G_a\] a canonical isomorphism, where \(\Lambda\) is the first etale cohomology of the spectral curve. The introduction reuses \(f\) inconsistently. Interpret the source using \(q:X\times\mathcal A\to\mathcal A\) and the target using \(\pi':Z\to\mathcal A\), exactly as the body does. The asserted isomorphism is false even under that well-typed reading.\par
\subsection*{Parameters, assumptions, and dimensions}
Take \(K=\mathbb C_3\), \(G=\operatorname{GL}_2\), and let \(X\) be the smooth projective completion of \(y^2=x^5-x\). The polynomial has five simple roots, since its greatest common divisor with its derivative is \(1\). Including the branch point at infinity, the hyperelliptic cover has six branch points and genus \(2\). Thus \(X\) is a smooth projective connected curve in characteristic zero and \(G\) is a permitted connected reductive group.\par
Choose \(b=(0,0)\) in the full Hitchin base. There is no regular-spectral-fibre restriction in Theorem 1.4.1. Use the monic characteristic polynomial explicitly displayed at the start of Section 5.5: the spectral equation is \(t^2=0\). The finite flat degree-two cover is the nilpotent thickening with \[\pi_*\mathcal O_{Z_b}=\mathcal O_X\oplus\Omega^{-1},
 \qquad \operatorname{Lie}J_b=\mathcal O_X\oplus\Omega^{-1}.\] The Hodge–Tate sequence used in (6.1.13) gives, at this fibre, \[\dim_K H_v^1(X,\mathcal O_X)=h^1(\mathcal O_X)+h^0(\Omega)=2+2=4.\] It also gives \[\dim_K H_v^1(X,\nu^*\Omega^{-1})
 =h^1(\Omega^{-1})+h^0(\mathcal O_X)=3+1=4,\] because Serre duality and Riemann–Roch give \(h^1(\Omega^{-1})=h^0(\omega_X^2)=3\). Tate twists do not change these dimensions. The source fibre therefore has dimension \(8\).\par
The etale site is unchanged by this nilpotent thickening, so \(H^1_{\mathrm{et}}(Z_b,\mathbb Z_p)=H^1_{\mathrm{et}}(X,\mathbb Z_p)\) has rank \(4\). Hence the target fibre has dimension \(4\), and \(\psi\) cannot be an isomorphism. More generally the same calculation gives dimensions \(5g-2\) and \(2g\) for \(g\geq2\). The nilpotent summand is invisible on reduced perfectoid test spaces.\par
\subsection*{Existing counterevidence and local correction}
Theorem 7.0.1 and Corollary 7.0.2, p. 30, use a comparison morphism and the Cartesian diagram. Remark 7.0.5 on that page explicitly warns that this map need not be an isomorphism, including for nonreduced spectral curves; source lines 1944–1945 corroborate this. Proposition 7.2.3, p. 33, establishes the needed kernel isomorphism. The Cartesian property is enough for the exponential splitting and the later homeomorphism argument.\par
The precise correction is to replace the introductory isomorphism claim by the comparison morphism and Cartesian assertion of Theorem 7.0.1/Corollary 7.0.2. The body already contains the necessary distinction. Consequently this is an individually Minor local overstatement: no invalid downstream derivation of the main correspondence has been demonstrated from it.\par
The displayed ideal formula near Definition 5.5.1 omits the monic term in its typography. The counterexample above deliberately uses the correct monic polynomial already given in the same subsection; it does not exploit that omission.\par
\section*{Problem 3: the dense smooth connected spectral locus needs a genus condition}
\subsection*{Claim and counterexamples}
Section 5.5, p. 22, source line 1411, states without a genus restriction that there is an open dense \(\mathcal A^\circ\subseteq\mathcal A\) over which the spectral family is smooth proper with geometrically connected fibres.\par
Take again \(K=\mathbb C_3\) and \(G=\operatorname{GL}_2\), but let \(X=\mathbb P^1\). These satisfy the subsection’s general curve and group assumptions. Since \(\Omega\simeq\mathcal O(-2)\) as a line bundle, \[H^0(X,\Omega)=H^0(X,\Omega^2)=0.\] Thus \(\mathcal A\) is the single zero point and its spectral cover has equation \(t^2=0\). It is a nonreduced finite flat double cover of \(X\), so it is not smooth. A dense open of this one-point base must contain that point; none has the asserted property.\par
For an elliptic curve, \(\Omega\) is trivial and the coefficients of the degree-two spectral polynomial are constants in \(K\). If it has distinct roots, the cover is the disjoint union of two copies of \(X\), hence is smooth but disconnected. A repeated root gives a nonreduced cover. Thus there is no smooth geometrically connected spectral fibre in this case either.\par
\subsection*{Impact and counterevidence}
Theorem 1.1.3, p. 2, Section 8.5, p. 41, and Theorem 8.6.2, p. 42, already impose genus at least \(2\) where the regular-locus geometry is developed. The main twisted correspondence of Section 6 uses all spectral covers and does not require this dense locus. No additional main-theorem failure has been established.\par
Restrict the assertion in Section 5.5 to \(g(X)\geq2\), as in the later treatment. Rank one can be handled separately: its spectral cover is \(X\) itself for every genus. This is Minor, a supported local scope defect with a matching later restriction.\par
The constant-coordinate description of \(\tau_\Omega\) in the same subsection should also be read with \(n\geq2\); in rank one the characteristic relation gives \(T=b\). Section 6.5 uses the appropriate rank-one dependence. This is retained as an unrated scope clarification, not an additional substantive finding.\par
\section*{Dependence, remaining work, and conclusion}
The repairs in Problems 1–3 are distinct: an analytic replacement proof, restoration of the body’s weaker comparison statement, and restoration of a later genus condition. None constitutes verification of an author revision.\par
The review did not independently prove all external inputs: the cited local \(p\)-adic Simpson results, regular-centralizer theory, rigid cohomological comparison, Lütkebohmert’s theorem, almost-mathematical input, algebraic/rigid GAGA antecedents, and Picard/BNR representability remain dependencies. The uniform etale local-constancy issue across all \(p\)-power levels in Lemma 8.6.1 and the full Ext/compactness argument in Lemma 8.6.8 were not independently resolved. The genus-zero scope of the openness argument in Section 9 also remains to be reconciled: intermediate statements 9.1.2 and 9.2.3 use \(g\geq1\), while 9.2.1 is more general. No counterexample or supported additional defect is claimed for these pending items. Ambiguities in regular-semisimple terminology at branch points were not promoted to findings.\par
The official arXiv record inspected during this review listed v1, submitted 2 February 2024. This was a bounded version check, not an exhaustive erratum search or a comparison with DOI journal-final content. No empirical experiment is needed for the three stated checks; the outstanding scientific work is the stated mathematical/dependency review.\par
Item | Judgment | Repair verification | Remaining status\par
1 | Major: core proof premise false | Exact lemma conclusion restored, conditional on 8.2.5 | Human adjudication pending\par
2 | Minor: introductory isomorphism false | Body’s morphism and Cartesian statement suffice | No new body defect shown\par
3 | Minor: genus scope omitted | Later \(g\geq2\) restriction restores stated scope | Rank one separate\par
Review | Bounded claim-focused audit | All 49 source pages visually checked | External dependencies pending\par
Build | Complete | Exact current TeX compiled | Changed/reflowed pages inspected\par
\section*{Consulted sources}
1. B. Heuer and D. Xu, p-Adic non-abelian Hodge theory for curves via moduli stacks, arXiv:2402.01365v1 (2024), supplied PDF and matching author TeX. Stable version link: https://arxiv.org/abs/2402.01365v1. DOI identity: https://doi.org/10.1090/jams/1065.\par
2. Official arXiv submission record, https://arxiv.org/abs/2402.01365, inspected during this task; archived evidence metadata identifies retrieval reference `turn20view0`.\par
3. Task-local source manifest, source excerpts, independent checks and page images identified in the appendix. External works cited by the manuscript were traced as dependencies rather than represented as fully reread sources.\par
\section*{Supporting checks}
All paths here are relative to this task’s assigned staging root.\par
- Original: `papers/0038/0038.pdf`; inventory: `papers/0038/source\_manifest.json`.\par
SHA-256:
  > b0c177126f1907eb426d2b9bd414da4fb
> 7edd0e644000ba70e325136e390d6cb\par
- Extracted full text SHA-256:
  > 3b585ed6d4fa3abdfd939a2c89f92afa7
> e0f82a3ec81173610b428004b49b2ea\par
It contains page text/images 001–049, extraction/render logs, `source\_search.json`, `source\_excerpts.json`, `web\_version.json`, and the independent-check script/result. The source excerpt ranges actually read were 481–504, 1371–1438, 1939–1952 and 2390–2426.\par
- The main assistant’s own 24-line `independent\_checks.py` was run using local Python, exit code 0. It checked finite geometric-series identities for \(1\leq N\leq8\), squarefreeness of \(x^5-x\), and the dimension arithmetic for \(2\leq g\leq6\). Infinite-series convergence, Laurent-unit classification, cohomological dimensions, the general repair and genus-boundary arguments are analytic/algebraic reasoning, not claims of finite computation proving an infinite theorem.\par
- Poppler extraction and rendering were performed by the task’s own `review\_io.py`. All 49 images were actually displayed and checked by the main assistant; render success alone was not counted as inspection. The first-page renderer emitted a Symbol display-font warning; the displayed page was readable. The visual-inspection ledger records the distinction between rendering and inspection.\par
No author code was executed. No author-source compilation, installation, extra quota, author contact or repository execution was performed.\par
\section*{Supporting checks}
The main assistant performed the scientific reasoning, static scientific-source inspection, independent calculations and image checks. A separate administrative agent was restricted to read-only metadata consistency, not scientific acceptance.\par
This assistant did not personally inspect the tab or its occupancy. Backend identity and equivalence of those names are unverified; no independent evidence of a backend mismatch was obtained.\par
The original compiler-absence probe preceded the authorized portable Tectonic follow-up. Its build and page inspections remain preserved as historical evidence. The current report is compiled from this exact TeX with Tectonic 0.17.0 in untrusted, cached-only mode. Exact input, dependency, log, PDF and page-image hashes, together with actual inspection of changed or reflowed pages, are recorded in the accompanying provenance. Original-review actions retain their original attribution; focused AI acceptance is recorded separately. Scientific adjudication remains pending, with no human review or author adoption claimed.\par
\end{document}
